%%vsgtu709

%

% %%%

%%%   Самарский государственный технический университет

%%%   Кафедра прикладной математики и информатики

%%%

%%%   Преамбула для статьи в журнал "Вестник Самарского государственного

%%%   технического университета. Серия: «Физико-математические науки»"

%%%   Ориентирована под MikTEx 2.4 for Win32

%%%

%%%   Хотя сам журнал собирается под GNU/Linux на дистрибутиве TeXLive



\documentclass[11pt,a4paper,twoside]{article}



\usepackage[cp1251]{inputenc}

\usepackage[T2A]{fontenc}

\usepackage[english,russian]{babel}

\usepackage{samgtubib}

\usepackage[dvips]{graphicx}

\usepackage[multiple]{footmisc}



\usepackage{samgtu}

\renewcommand{\VestnikYear}{2009} % Год издания Вестника

\renewcommand{\VestnikNumber}{2\,(19)} % Номер Вестника

\renewcommand{\VestnikMonth}{Ноябрь} % Месяц выхода Вестника

\renewcommand{\VestnikData}{01.11.09} % Дата подписания Вестника в печать

\renewcommand{\VestnikUPL}{26,64} % Усл. печ. л.

\renewcommand{\VestnikUIL}{25,73} % Уч.-изд. л.

\renewcommand{\VestnikRegNumber}{479/09} % Регистрационный номер

\renewcommand{\VestnikOrder}{994} % Номер заказа







%%

% \begin{document}





\renewcommand{\firstpage}{144}

\renewcommand{\lastpage}{151}

% \Partition{Математическое моделирование}{Applied mechanics} {Applied mechanics} % Раздел Вестника, в который подаётся статья

% Названия разделов см. на сайте при подаче заявки





%%%%%%%%%%%%%%%%%%%%%%%%%%%%%%%%%%%%%%%%%%%%%%%%%%%%%%%%%%%%%%%%%%%%%%%%%%%%%%%%%%%%

%Информация о статье и об авторах на русском и английском языках



% Номер УДК

\udk{519.864.3}

\msc{91B30; 62P05}



\title{Рекуррентные методы построения кумулятивных функций риска}% Полный заголовок

{Рекуррентные методы построения кумулятивных функций риска }% Заголовок, который идёт в верхний колонтитул



\author{В.\,Н.~Никишов, Е.\,В.~Михайлова }% Автор(ы) для заголовка, если авторов несколько укать \inst

{Н\,и\,к\,и\,ш\,о\,в~В.\,Н., М\,и\,х\,а\,й\,л\,о\,в\,а~Е.\,В.}% Автор(ы) для оглавления и колонтитула через \,



\institute{\samgu.}

% Если несколько, то так  \institute{ Организация 1 \and Организация 2  \and Организация 3}



\email{E-mails}{tsh-sea05@yandex.ru, milena82@yandex.ru}  %E-mail's авторов



\otherinf{\fullauthor{Виктор Николаевич Никишов}\regalia{к.ф.-м.н., доц.} \position{доцент}\dept{каф. математики и бизнес"=информатики}

\fullauthor{Елена Владимировна Михайлова} \position{аспирант} \dept{каф. математики и бизнес"=информатики}  }



\keywords{суммарный размер убытка, рекуррентные методы, кумулятивные функции риска}



\workabstract{Рассматривается задача оценки суммарных убытков предприятия, связанных с~ненадлежащим выполнением обязательств. 

Функция и~плотность распределения суммарного размера убытков для совокупности договоров строятся по рекуррентным методам, которые были предложены  N.~de Pril и H.~Panjer. 

Выполнены численные эксперименты на основе пяти групп портфелей договоров. 

Проведён анализ результатов, который позволил установить достоинства и~недостатки алгоритмов расчёта. 

В~частности, значения средств под риском, полученные по методу H.~Panjer, превышают соответствующее значение, полученное по методу N.~de Pril, а максимальное значение вероятности наступления убытка в~функции плотности распределения, полученное по методу H.~Panjer, занижено по сравнению со значением, полученным по методу N.~de Pril.

Результаты могут быть использованы для дальнейшего развития подходов N.~de Pril и H.~Panjer к оценке суммарного размера убытков совокупности рисков.}



\engtitle{Recurrent methods of construction of cumulative functions of risk}



\engauthor{V.\,N.~Nikishov, E.\,V.~Mikhailova}

{N\,i\,k\,i\,s\,h\,o\,v~V.\,N., M\,i\,k\,h\,a\,i\,l\,o\,v\,a~E.\,V.}





\enginstitute{\engsamgu.}% Место работы каждого из авторов на англ. языке



\otherenginfm{\fullauthor{Viktor N. Nikishov} \regalia{Ph.\,D. (Phys. \& Math.)} \position{Associate Professor} \dept{Dept. of Mathematics and Business Informatics}

\fullauthor{Elena V. Mikhailova} \position{Postgraduate Student} \dept{Dept. of Mathematics and Business Informatics}}



\engabstract{The estimation problem of the enterprise total losses related to the improper performance of liability is considered. The total losses distribution function and density for the aggregate contracts are constructed by the recurrent methods, proposed by N.~de Pril and H.~Panjer. The numerical experiments based on five groups of contract portfolio are carried out. The results analysis showing the advantages and disadvantages of calculating algorithms is made. In particular, the values of means at risk obtained by H.~Panjer method are superior to the corresponding values founded  by N.~de Pril method, and the maximum value of loss occurrence probability in the distribution density realized by H.~Panjer method is less than the one, gotten by N.~de Pril method. The results can be used for the further development of N.~de Pril and H.~Panjer approaches for estimation of total losses of risk aggregate.}



\engkeywords{total loss, recurrent methods, cumulative functions of risk}



\date{24/XI/2011}

\revisiondate{27/III/2012}



%%%%%%%%%%%%%%%%%%%%%%%%%%%%%%%%%%%%%%%%%%%%%%%%%%%%%%%%%%%%%%%%%%%%%%%%%%%%%%%%%%%%





\maketitle





\Section[n]{Введение} 

Осуществление предпринимательской деятельности неразрывно связано с~заключением договоров, обязательства по которым не всегда исполняются надлежащим образом и в срок. 

Таким образом, актуальной для предпринимателя является задача оценки суммарного размера убытков.



Рассмотрим производственное предприятие, которое поставляет продукцию потребителям  и/или закупает материалы для производства на основе коммерческих контрактов. 

Совокупность из $N$ контрактов предприятия  можно разделить на группы исходя из вероятности наступления убытков и их возможного размера. 

В~настоящее время для оценки риска на предприятии рассматривается несколько групп риска (от минимального до катастрофического) \cite{Mihailova:bib:Granaturov}. Будем считать, что количество групп равно $K$.



Пусть $X(n)$ "--- случайная величина, описывающая возможный размер убытка предприятия по $n$-ному ($n=1, 2, \ldots N$) договору совокупности, может быть представлена как

\begin{equation}

\label{nik:eq:1}

 X(n)=I(n) \, Y(n),

\end{equation}

где $I(n)$ "--- индикатор, который отражает факт наступления события, связанного с получением ущерба по договору с~$n$-ным контрагентом, и подчиняется закону распределения Бернулли (величина $I(n)$ принимает значения 1 или 0 с вероятностями $q(n)$ и $p(n)=1-q(n)$ соответственно); $Y(n)$ "--- дискретная случайная величина, которая принимает $J$ значений, описывает фактический размер ущерба в случае, когда $I(n)=1$, и подчиняется некоторому закону распределения $f(l)$, например, бета-распределению. Значения $Y(n)$ фактического размера убытка нормированы, т.\,е. могут быть представлены выражением

\begin{equation}

\label{nik:eq:2}

Y(n)=l (n) \, B, 

\end{equation}

где $B$ "--- порядок убытков, $l(n)$ "--- нормированное значение $Y(n)$. Далее порядок убытков $B$ будем считать равным единице.



Для вычисления функции и плотности распределения совокупного размера убытков совокупность контрактов представим в~виде таблицы, в каждой строке $k=1, 2, \ldots, K$ которой указывается вероятность неисполнения обязательств по договорам $q_k$, 

а в каждом столбце  $j=1, 2,\ldots, J$ "--- возможный размер убытка $l_j$, в~ячейке "--- количество договоров  $N_{kj}$, вероятность неисполнения обязательств и возможный размер убытков по которым равен $q_{k}$ и  $l_{j}$, соответственно.

С учётом разбиения совокупности договоров на ячейки выражения \eqref{nik:eq:1}, \eqref{nik:eq:2} примут вид

$$

X_{kj}(n)= I_{kj}(n) \,Y_j (n),\quad  Y_j(n)=l_j(n) \, B.

$$





Суммарный убыток $\widetilde{S}$ предприятия является случайной величиной и принимает значения $\widetilde{S} = 0, 1, 2, \dots, [S_{\max}]$, где 

\begin{equation}

 \label{f1.2}

S_{\max} =

\sum_{k=1}^K \sum_{j=1}^J \sum_{n=1}^{N_{kj}} X_{kj}(n)= 

\sum_{k=1}^K \sum_{j=1}^J \sum_{n=1}^{N_{kj}} l_{kj}(n) = 

\sum_{k=1}^K \sum_{j=1}^J l_{kj} \,  N_{kj}.

\end{equation}

Здесь $l_{kj}(n)$ "--- нормированное значение фактического размера убытка $Y_j(n)$ в случае, когда $I_k(n)=1$;

$l_{kj}$  "--- нормированное значение фактического размера убытка $Y_j$ в случае, когда вероятность наступления убытка равна $q_k$. В~последнем равенстве формулы \eqref{f1.2} осуществляется переход от возможного убытка предприятия от $n$-ного договора к фактическому убытку в ячейке~$N_{kj}$.



Таким образом, предприятие может получить суммарный убыток $\widetilde{S}$ с вероятностью

\begin{equation}

\label{f.2}

p(s)=P( s\geq \widetilde{S})=\sum_{s=0}^{\widetilde{S}} p(s).

\end{equation}

База рекурсии для формулы \eqref{f.2} будет определена ниже с~помощью рекуррентных методов N.~de~Pril и H.~Panjer, основанных на модели индивидуального и~коллективного рисков соответственно.



Получение необходимых выражений для вычисления вероятности $p(s)$ того, что убыток примет значение $\widetilde{S}$, основано на использовании производящей функции, которую можно записать в виде  \cite{Mihailova:bib:Sevastianov}

\begin{equation}

\label{f3.2}

H(t)=\sum_{s=0}^{\widetilde{S}}p(s)\, t^s,

\end{equation}

где $t$ "--- вспомогательная переменная.



\smallskip



\Section[n]{Рекуррентный метод N.~de~Pril} 

Рекуррентный метод N.~de Pril \cite{Mihailova:bib:De_Pril1, Mihailova:bib:De_Pril2} основан на индивидуальной модели риска, согласно которой возникающие убытки определяются контрактом и~каждый контракт может привести к~убытку один раз \cite{Mihailova:bib:Falin, Mihailova:bib:Korolev}.



Выражение для расчёта вероятности $p(s)$ того, что убыток примет значение $\widetilde{S}$, получено методом N.~de~Pril  на основе производящей функции суммарного размера убытков:

\begin{equation}

\label{f.2'}

H(t)=\prod_{k=1}^{K}\prod_{j=1}^{J}(p_k+q_k t^j) ^{N_{kj}},\quad p_k=1-q_k.

\end{equation}

Найдём производную функции \eqref{f.2'} по $t$, используя её логарифмирование:

\begin{equation}

\label{f3.10}

\frac{H'(t)}{H(t)}=

\sum_{k=1}^K\sum_{j=1}^J N_{kj} \, q_k \, j \, t^{j-1} \, M, 

\end{equation}

где $M = (p_i+q_i t^j)^{-1}$.

Используя разложение в ряд множителя $M$, выражение \eqref{f3.10} преобразуем к виду

\begin{equation}

\label{f.5}

\frac{H'(t)}{H(t)}=

\sum_{j=1}^J \sum_{l=1}^{\infty} \biggl( t^{jl-1} \, j \, \sum_{k=1}^K (-1)^{l-1} \, N_{kj} \Bigl (\frac{q_k}{p_k}\Bigr)^l \biggr).

\end{equation}

Обозначив через

\begin{equation}

\label{f3.11}

A(l-1,j)=j\, \sum_{k=1}^K (-1)^{l-1} \, N_{kj} \Bigl (\frac{q_i}{p_i}\Bigr)^l,

\end{equation}

выражению \eqref{f.5} придадим вид

\begin{equation}

\label{f3.12}

H'(t)=H(t)\cdot \sum_{j=1}^{J}\sum_{l=1}^{\infty} t^{jl-1} \, A(l-1, j).

\end{equation}



Подставляя в \eqref{f3.12} представление $H(t)$ в виде \eqref{f3.2}, получим

\begin{equation}

\label{f3.13}

H'(t)=\sum_{j=1}^{J}\sum_{l=1}^{\infty}\sum_{s=0}^{\widetilde{S}}t^{jl+s-1} \, A(l-1, j) \, p(s).

\end{equation}

Изменим выражение \eqref{f3.13}, используя производную производящей функции~\eqref{f3.2}:

\begin{equation}

\label{f3.20}

\sum_{s=1}^{\widetilde{S}}p(s) \, s \, t^{s-1}=

\sum_{j=1}^{J}\sum_{l=0}^{\infty}\sum_{s=0}^{\widetilde{S}}t^{jl+s-1} \, A(l-1, j)\, p(s).

\end{equation}

Приравнивая в \eqref{f3.20} множители при одинаковых степенях $t$, %и учитывая, что в правой части $s$ "--- индекс суммирования, 

получим выражение для расчёта искомой вероятности $p(s)$:

\begin{equation}

\label{f3.15}

p(s)=\frac{1}{s}\sum_{j=1}^{\min\{s, J\}}\sum_{l=0}^{[{s}/{j}]} A(l-1, j) \, p(s-j l),

\end{equation}

где $s = 1, 2, \dots, [S_{\max}]$



Начальное значение (база рекурсии) для расчёта по \eqref{f3.15} имеет такой вид:

\begin{equation}

\label{f3.16}

p(0)=p_0\equiv\prod_{k=1}^{K} (1-q_k)^{\sum_{k=1}^{K}N_{kj}}.

\end{equation}



\smallskip



 \Section[n]{Рекуррентный метод H.~Panjer} 

 Получим выражение для расчёта вероятности $p(s)$ того, что убыток примет значение $s$, методом H.~Panjer~\cite{Mihailova:bib:Panjer}, который основан на коллективной модели риска. 

 В~этом случае ущерб определяется выражением

 \begin{equation}

 \label{f2.2}

S= Y_{1}+ Y_{2}+\ldots + Y_{\nu},

\end{equation}

где $\nu$ "--- случайная величина, которая описывает количество возможных исков по портфелю контрактов в~целом и имеет закон распределения $$\pi(\widetilde{\nu})=P(\nu=\widetilde{\nu}),$$ $\widetilde{\nu}$ "--- число убытков.



В работе \cite{Mihailova:bib:Panjer} доказано, что рекуррентная процедура расчёта суммарного размера ущерба применима для любых распределений, удовлетворяющих 

\begin{equation}

\label{f4.6}

\pi(\widetilde{\nu})=\Bigl(a+\frac{b}{\widetilde{\nu}}\Bigr) \, \pi(\widetilde{\nu}-1),\quad  \widetilde{\nu}\geqslant 1,

\end{equation} 

где $a$, $b$ "--- константы, определяемые процессом поступления исков.



Записывая производящую функцию для суммарного размера убытка $H_S(t)$ с~учётом закона распределения числа убытков $\pi (k)$ и

производящей функции возможного размера ущерба $H_Y(t)$, получим

\begin{equation}

\label{f4.4}

H_S(t)=\pi(0)+\sum_{k=1}^\infty\pi(k)\, [H_Y(t)]^k.

\end{equation}

С учётом \eqref{f4.6} после ряда преобразований производная функции \eqref{f4.4} примет вид

\begin{equation}

\label{f4.7}

H'_S(t)=(a+b)\, H'_Y(t) \, H_S(t)+a \, H_Y(t) \, H'_S(t).

\end{equation}

Записывая \eqref{f4.7} в развернутом виде с учётом производящих функций возможного размера убытка и функции \eqref{f3.2}, получим

\begin{equation}

\label{f4.10}

H'_S(t)=(a+b)\, \sum_{j=1}^\infty \sum_{s=0}^\infty j \, p(j) \, P_S(s) \, t^{s+j-1}+a \, \sum_{j=0}^\infty \sum_{s=1}^\infty 

s \, p(j) \, P_S(s) \, t^{s+j-1}.

\end{equation}

Здесь величина $P_S(s)$ означает вероятность наступления суммарного размера убытка $s$ и по сути является величиной $p(s)$, введённой по формуле \eqref{f.2};  $p(j)$ "--- вероятность наступления фактического размера убытка $Y=j$.



Приравнивая правые части \eqref{f4.10} и производной производящей функции \eqref{f3.2}, получим

\begin{multline}

\label{f4.11}

\sum_{s=1}^{\widetilde S} P_S(s) \, s \, t^{s-1} =

(a+b) \, \sum_{s=0}^{\widetilde S} \sum_{j=1}^\infty j \, p(j) \, P_S(s) \, t^{s+j-1}+\\

+ a\, \sum_{s=1}^{\widetilde S} \sum_{j=0}^\infty s \,p(j) \,P_S(s)\, t^{s+j-1}.

\end{multline}

После преобразования с учётом постановки задачи  \eqref{f4.11} примет вид

\begin{equation}

\label{f4.12}

 \sum_{s=1}^{\widetilde S} P_S(s)\, s\, t^{s-1} = s\, a\, p(0)\, P_S(s)

 + \sum_{s=1}^{\widetilde S} \sum_{j=1}^{\min\{s, J\}} p(j)\, P_S(s-j)\, (as+bj) \,t^{s-1}.

\end{equation}

Приравнивая в \eqref{f4.12} коэффициенты при одинаковых степенях $t$, получим

\begin{equation}

\label{f4.13}

P_S(s) = \frac{1}{1-a\, p(0)} \sum_{j=1}^{\min\{s, J\}} p(j) \, P_S(s-j)\, \bigl(a+{b\,j}/{s}\bigr).

\end{equation}



Начальное условие (база рекурсии) для расчёта по формуле \eqref{f4.13}) получим,  подставляя в \eqref{f4.4} $t=0$ с учётом производящей функции для возможного размера убытков:

$$P_S(0) =H_S(0)=\pi(0) + \sum_{k=1}^\infty \pi(k) \, [p(0)]^k.$$

Так как  $p(0)=p(Y=0)=0$, \eqref{f4.13} можно записать в виде

\begin{equation}

\label{f4.14}

P_S(s)= \sum_{j=1}^{\min\{s,J\}} p(j)\, P_S(s-j)\, \bigl(a+{b\, j}/{s}\bigr)

\end{equation}

с начальным условием

\begin{equation}

\label{f4.15}

P_S(0)= H_S(0) = \pi(0)=P(\nu=0).

\end{equation}



Для закона распределения Пуассона формулы \eqref{f4.14} и \eqref{f4.15} запишутся так:

\begin{equation}

\label{f4.17}

P_S(n)=\sum_{k=1}^{\infty} p(k) \, P_S(n-k) = \lambda \, k/{n}, \quad P_S(0)=\pi(0) = \exp({-\lambda}).

\end{equation}





\smallskip

\Section[n]{Анализ численных результатов} По формулам \eqref{f3.15} и \eqref{f3.16} для метода N.~de Pril и \eqref{f4.17} для метода H.~Panjer для модельного примера  были получены функции плотности распределения  суммарного размера убытков $f(s)$ совокупности рисков при $N=300$ с~максимально возможным размером убытков $S_{\max}=1636$~усл.~ед., графики которых представлены на рисунке слева.



\begin{figure}[h] \centering

\includegraphics[width=6cm]{./00PIC/graf1} \hfil \includegraphics[width=6cm]{./00PIC/graf2} \\

\caption*{Функции плотности распределения суммарного размера убытков $f(s)$ и распределения суммарного размера убытков $F(s)$ для модельного примера: \scriptsize сплошная линия "--- метод N.~de Pril; штриховая линия "--- метод H.~Panjer}



\end{figure}



Далее, используя значения функции плотности распределения, для рассматриваемого модельного примера были получены функции распределения суммарного размера убытков $F(s)$, графики которых представлены на рисунке справа.







Полученная функция распределения суммарного размера убытков позволяет получить значение величины средств под риском. 

Средства под риском  (Var) "--- выраженная в~заданных условных единицах оценка ожидаемых в~течение данного периода времени и с~заданной вероятностью потерь под воздействием различных факторов риска. 

В данной работе значение Var определялось с~вероятностью~0,95.



На основе численного эксперимента, выполненного на модельных примерах, были получены результаты, которые обобщены в~таблицу.



В строках таблицы представлены результаты по каждой совокупности. 

В~первом столбце таблицы указывается количество договоров $N$ в совокупности, 

во втором "--- максимальный размер возможного размера убытка по совокупности $S_{\max}$ (в~усл. ед.). 

В третьем и шестом  "--- время работы программ, реализующих методы; 

в четвёртом и седьмом "--- порядок количества итераций для рассматриваемых методов. 

В пятом столбце приведены значения средств под риском (в~усл. ед.), иллюстрирующие оценку максимальных потерь портфеля предпринимательских договоров с~вероятностью $0,95$, рассчитанных по функции распределения суммарного размера убытков по методу N.~de Pril. 

В~восьмом столбце приводятся отклонения значений средств под риском от значений, приведенных в пятом столбце, для метода H.~Panjer.



\begin{table}[b!]

\centering

\small

\begin{tabular}{r|r|rrr|r|r|r|r|l}

\hline

\multicolumn{1}{c|}{\footnotesize $N$} &  

\multicolumn{1}{c|}{\footnotesize $S_{\max}$}  &

\multicolumn{5}{c|}{\footnotesize метод N.~de Pril} & 

\multicolumn{3}{c}{\footnotesize метод H.~Panjer}  \\

\hline

\multicolumn{1}{c|}{\footnotesize 1 } & \multicolumn{1}{c|}{\footnotesize 2}   & 

\multicolumn{3}{c|}{\footnotesize 3} & \multicolumn{1}{c|}{\footnotesize 4}   & 

\multicolumn{1}{c|}{\footnotesize 5} &  

\multicolumn{1}{c|}{\footnotesize 6} & 

\multicolumn{1}{c|}{\footnotesize 7} & 

\multicolumn{1}{c}{\footnotesize 8}  \\

\hline

\rule{0mm}{12pt}

  150  &    818   &      &        &  1 917 мс & 7  &   281     & $< 1$ мс  &  5 & $+ 15$       \\

  300  &  1 636   &      &        &  7 381 мс & 8  &   531     & $< 1$ мс  &  5 & $+ 21$   \\

  750  &  4 090   &      &        & 45 205 мс & 9  &  1 260    & $< 1$ мс  &  5 & $+ 33$    \\

1 500  &  8 180   &      &  2 мин & 45 758 мс & 9  &  2 451    & $< 1$ мс  &  6 & $+ 46$   \\

3 000  & 16 360   &      &  8 мин & 18 845 мс & 10 &  4 806    & $< 1$ мс  &  6 & $+ 64$   \\

7 500  & 40 900   & 1 ч  &  1 мин & 37 724 мс & 11 & 11 800    & $< 1$ мс  &  6 & $+ 100$   \\

15 000 & 81 800   & 2 ч  & 58 мин & 8 768 мс  & 11 & 23 383    & $< 1$ мс  &  6 & $+141$  \\

\hline

\end{tabular} 



\end{table}



Например, первая строка таблицы  содержит результаты численного эксперимента для совокупности с~количеством договоров $N=150$ и максимально возможным размером ущерба $S_{\max}=818$~усл.~ед. 

При этих условиях средства под риском  по методу N.~de Pril  составили 281~усл.~ед, а по методу H.~Panjer "--- на 15~усл.~ед. больше, т.\,е. 296~усл.~ед.

Программа на основе метода N.~de Pril выполнялась 1~917~мс и~выполнила порядка $10^7$ итераций,

а программа на основе метода H.~Panjer "--- менее 1~мс и порядка $10^5$ итераций.



Рассмотренные методы позволяют заключить, что время работы и~количество итераций алгоритмов зависят только от величины $S_{\max}$.

На основе анализа сводной таблицы и полученных графиков можно сделать следующие выводы:

\begin{itemize}

\item[1)] установлено влияние диапазона вероятностей и~групп риска в совокупности на величину $S_{\max}$;

\item[2)] несмотря на одинаковое поведение функций и~плотностей распределения суммарного риска, значение величины средств под риском изменяется в~зависимости от диапазона вероятностей в совокупности.

\end{itemize}



Особенность метода H.~Panjer, состоящая в том, что один контракт может привести более чем к~одному иску, отражается в полученных результатах следующим образом:

\begin{itemize}

\item[1)] значения средств под риском, полученные методом H.~Panjer, превышают соответствующее значение, полученное методом N.~de Pril;

\item[2)] максимальное значение вероятности наступления убытка в~функции плотности распределения, полученной по методу H.~Panjer, занижено по сравнению с методом N.~de Pril.

\end{itemize}



Полученные результаты могут быть использованы для дальнейшего развития подходов N.~de Pril и H.~Panjer к оценке суммарного размера убытков совокупности рисков.



\begin{thebibliography}{9}



\RBibitem{Mihailova:bib:Granaturov}

\by Гранатуров~В.~М.

\book  Экономический риск: сущность, методы измерерния, пути снижения

\publaddr М.

\publ Дело и сервис

\yr 1999

\totalpages 111

\misctransl

\by Granaturov~V.\,M.

\book Economic risk. Essence, Methods of Measuring, Ways of Reduction

\publaddr Moscow

\publ Delo i Servis

\yr 1999

\totalpages 111



\RBibitem{Mihailova:bib:Sevastianov}

\by Севастьянов~Б.~А.

\book  Курс теории вероятностей и математической статистики

\publaddr М.

\publ Наука

\yr 1982

\totalpages 256

\misctransl

\by Sevastyanov~B.\,A.

\book A~Course in Probability Theory and Mathematical Statistics

\publ Nauka

\publaddr Moscow 

\yr 1982

\totalpages 256





\Bibitem{Mihailova:bib:De_Pril1}

\paper On the exact computation of the aggregate claims distribution in the individual life model

\by De~Pril~N.

\jour Astin Bulletin

\yr 1986

\vol 16

\issue 2

\pages 109--112



\RBibitem{Mihailova:bib:De_Pril2}

\paper The aggregate claims distribution in the individual model with arbitrary positive claims

\by De~Pril~N.

\jour Astin Bulletin

\yr 1989

\vol 19

\issue 1

\pages 9--24



\RBibitem{Mihailova:bib:Falin}

\by Фалин~Г.~И.

\book  Математический анализ рисков в страховании

\publaddr Москва

\publ Российский юридический издательский дом

\yr 1994

\totalpages 130

\misctransl

\book Mathematical Analysis of Risks in Insurance

\by Falin~G.\,I.

\yr 1994

\totalpages 130

\publ Rossiyskiy Yuridicheskiy Izdatel'skiy Dom 

\publaddr Moscow



\RBibitem{Mihailova:bib:Korolev}

\by Королев~В.~Ю., Бенинг~В.~Е., Шоргин~С.~Я.

\book  Математические основы теории риска

\publaddr М.

\publ Физматлит

\yr 2007

\totalpages  544

\misctransl

\by Korolev~V.\,Yu., Bening~V.\,E., Shorgin~S.\,Ya.

\book Mathematical Foundation of Risk Theory

\publaddr Moscow

\publ Fismatlit

\yr 2007

\totalpages  544



\Bibitem{Mihailova:bib:Panjer}

\paper Recursive evaluation of a famaly of compound distributions

\by   Panjer~H.\,H.

\jour Astin Bulletin

\yr 1981

\vol 12

\issue 1

\pages 12--26







\end{thebibliography}



\noindent

\hskip6.5cm \thedate \par\noindent

\hskip6.5cm\therevdate



\pagebreak



\chead[\fancyplain{}{\scriptsize \sl N\,i\,k\,i\,s\,h\,o\,v~V.\,N., M\,i\,h\,a\,i\,l\,o\,v\,a~E.\,V.}]{\fancyplain{}{\scriptsize \sl  Recurrent methods of construction of cumulative functions of risk \dots}}



\makeengtitlem



\newpage







% \newpage

%

% \tableofengcontents

%

%\end{document}

